%!TEX root = ../main.tex

\chapter{Introduction\label{chap:introduction}}

\acp{TCG} such as Magic: The Gathering, Pokémon, and various sports card collections have evolved from 
simple hobbies into a global market with significant economic value. For collectors and local game stores, managing 
inventory is a critical task. This involves sorting thousands of cards, assessing their condition, and identifying their 
current market price. However, as collections grow, manual sorting becomes an increasingly time-consuming, tedious, and 
error-prone process.

While the demand for automation in this field is high, the current landscape of solutions is polarized. On one hand, 
there are industrial-grade sorting machines used by large companies. On the other hand, commercial solutions available 
to small businesses or individual collectors are often prohibitively expensive, costing thousands of dollars. 
Furthermore, these proprietary systems are typically \enquote{black boxes}, which are closed-source and almost 
impossible to modify or improve by the user.

This thesis addresses the gap between high-end commercial sorters and manual labor. The primary motivation for this work 
is the lack of accessible, affordable, and transparent tools for card automation. Existing consumer-grade machines often 
fail to meet specific requirements, are unavailable in certain regions, or lack the robustness required for handling 
valuable collectibles.

The goal of this project is to develop an automated system for recognizing and classifying trading cards that adheres to 
the philosophy of Open-source software and hardware. By designing a robot based on accessible components, such as an 
Arduino microcontroller and stepper motors, and combining it with modern computer vision techniques, this work aims to 
democratize the technology required for card sorting.

Beyond the practical application, this project serves an educational purpose. It demonstrates how complex concepts in 
robotics and machine learning can be applied to a relatable, real-world hobby. By making the design replicable, this 
thesis invites enthusiasts to build, operate, and further improve the platform at home.

\section{Related Works}

The problem of automated trading card processing involves two main domains: mechanical handling and software 
recognition. Existing solutions can be broadly categorized into commercial proprietary systems and academic or hobbyist 
research projects.

\subsection{Commercial Solutions}

Several commercial automated sorters exist on the market, primarily targeting card shops and high-volume sellers.

Systems such as \textit{CardBot} \cite{cardbot} and \textit{Card Dealer Pro} \cite{carddealerpro} represent the high-end 
segment. \textit{CardBot} utilizes silicone suction cups for gentle handling, which is a crucial feature for valuable 
collectibles. Similarly, \textit{TCG Machines} \cite{tcgmachines} claims advanced features like foil detection using 
specialized LED arrays and diffusers to overcome the reflective nature of holographic cards.

\subsection{Academic and Research Approaches}

In the research domain, several methodologies have been proposed for object detection and card classification.

\subsubsection{Computer Vision Techniques}
The use of \ac{YOLO} is prevalent in recent literature due to its real-time performance.
\begin{itemize}
    \item \textit{Card Detection:} A study on self-playing poker robots \cite{pokerrobot} successfully utilized 
    \ac{YOLO}v5 combined with a robotic arm using suction mechanisms, validating the feasibility of combining vision 
    with pneumatics.
    \item \textit{Verification:} Research in ID card verification \cite{idcardverification} demonstrates a pipeline 
    similar to trading cards: segmentation using \ac{YOLO}v8 followed by classification using \ac{ResNet}.
    \item \textit{Synthetic Data:} A significant challenge in training neural networks is the lack of labeled data. 
    Research on Pokémon card detection \cite{pokemonstanford} addressed this problem by creating a synthetic dataset, 
    i.e., artificially generating card images with simulated lighting and occlusions. This approach allowed them to 
    train a robust \ac{YOLO}v8 model without the need to manually photohraph thousands of cards.
\end{itemize}

\subsubsection{Siamese Networks}
An alternative to standard \ac{OCR} is visual similarity matching. Lowhur \cite{lowhuryugioh} demonstrated a system for 
Yu-Gi-Oh! cards achieving 99\% accuracy. Instead of reading text, the system uses Siamese networks and triplet loss to 
learn a \enquote{fingerprint} of the card's artwork and match it against a database. This highlights that for cards with 
distinct artwork, visual matching can be superior to text reading.

%%%%%% FROM MRS TEMPLATE:

\section{Contributions}

This section should describe the author's contributions to the field of research.

\section{Mathematical notation}

It is a good practice to define basic mathematical notation in the introduction.
See \reftab{tab:mathematical_notation} for an example.

\begin{table*}[!h]
  \scriptsize
  \centering
  \noindent\rule{\textwidth}{0.5pt}
  \begin{tabular}{lll}
    $\mathbf{x}$, $\bm{\alpha}$ & vector, pseudo-vector, or tuple\\
    $\mathbf{\hat{x}}$, $\bm{\hat{\omega}}$& unit vector or unit pseudo-vector\\
    $\mathbf{\hat{e}}_1, \mathbf{\hat{e}}_2, \mathbf{\hat{e}}_3$ & elements of the \emph{standard basis} \\
    $\mathbf{X}, \bm{\Omega}$ & matrix \\
    $\mathbf{I}$ & identity matrix \\
    $x = \mathbf{a}^\intercal\mathbf{b}$ & inner product of $\mathbf{a}$, $\mathbf{b}$ $\in \mathbb{R}^3$\\
    $\mathbf{x} = \mathbf{a}\times\mathbf{b}$ & cross product of $\mathbf{a}$, $\mathbf{b}$ $\in \mathbb{R}^3$\\
    $\mathbf{x} = \mathbf{a}\circ\mathbf{b}$ & element-wise product of $\mathbf{a}$, $\mathbf{b}$ $\in \mathbb{R}^3$ \\
    $\mathbf{x}_{(n)}$ = $\mathbf{x}^\intercal\mathbf{\hat{e}}_n$ & $\mathrm{n}^{\mathrm{th}}$ vector element (row), $\mathbf{x}, \mathbf{e} \in \mathbb{R}^3$\\
    $\mathbf{X}_{(a,b)}$ & matrix element, (row, column)\\
    $x_{d}$ & $x_d$ is \emph{desired}, a reference\\
    $\dot{x}, \ddot{x}, \dot{\ddot{x}}$, $\ddot{\ddot{x}}$ & ${1^{\mathrm{st}}}$, ${2^{\mathrm{nd}}}$, ${3^{\mathrm{rd}}}$, and ${4^{\mathrm{th}}}$ time derivative of $x$\\
    $x_{[n]}$ & $x$ at the sample $n$ \\
    $\mathbf{A}, \mathbf{B}, \mathbf{x}$ & LTI system matrix, input matrix and input vector\\
    \emph{SO(3)} & 3D special orthogonal group of rotations\\
    \emph{SE(3)} & \emph{SO(3)}~$\times~\mathbb{R}^3$, special Euclidean group\\
  \end{tabular}
  \noindent\rule{\textwidth}{0.5pt}
  \caption{Mathematical notation, nomenclature and notable symbols.}
  \label{tab:mathematical_notation}
\end{table*}
